\chapter*{Introduction}
\addcontentsline{toc}{chapter}{Introduction}

\section*{What this is}

This collection is a broad overview and survey. It walks through what the esoteric groups of the world have
been saying for thousands of years, in their long effort to model the universe and our experience of living in
this reality we call home. It does not take a side on whether any of it is literally true. It treats each
tradition as a serious attempt to map the same territory, the nature of reality and the nature of the self,
and it asks what they saw.

The word \textbf{esoteric} here means the inner, deeper layer of a tradition, the part concerned with the
hidden structure of reality and the inner transformation of the person, as opposed to the outer layer of
ritual and rule. Every tradition covered here has such an inner layer, and that is what we follow.

\section*{Why gather all of this}

There is a practical reason behind the contemplative one. To model consciousness, the inner life, the sense
of the divine, the path of growth, and the worlds of beings both positive and negative, it helps to know the
full depth and scope of what was already mapped. The traditions spent millennia naming the parts of the soul,
charting the levels of reality, describing the stages of inner ascent, and cataloguing the motifs that recur
in dream, vision, and myth. Before building any new model of mind or spirit, it is worth seeing the whole
inherited record, so that nothing important is left out and the new work can build on the old depth.

\section*{What we are looking for}

The first aim is the simplest. To see how each tradition actually \textbf{mapped reality}, the source, the
worlds, the make-up of the soul, the shape of time, and the path of a life. Each is a serious attempt to chart
the same territory, and these maps are the heart of the matter, gathered here so their full depth and scope
are visible at once.

The second aim is to look for the \textbf{thing behind the thing}. Underneath the many models there seems to
be one reality they are all pointing at, partly seen and reported in different words. Read together, the
traditions locate it from many directions, each from its own angle, its own partial view, its own depth of
experience. Naming that shared reality, the thing the models only point toward, is the main goal.

The differences matter less. No map is ever complete, and the variations are mostly different framings of the
one thing. Now and then a difference widens the picture or offers an alternative worth keeping, so we note
them, but lightly, as hints rather than as the point. Dwelling on them tends to distract from what the
traditions are together circling.

Stated in advance and in the plainest words, that shared shape looks like this. There is one source. The one
goes out into the many. The many forget the one. And the point of a life, and of a path, is to remember, to
return, to experience reality to the fullest, to tune back into the oneness, and to find the way home. This is
the flow behind the field, the way that lies beyond all the senses. Whether a tradition calls the return
liberation, union, repair, awakening, or salvation, the shape is strikingly constant. We let the maps build
that case rather than assume it.

\section*{How the collection is organized}

The traditions are grouped by how deep and explicit their inner modeling went, not by how true or important
they are.

\begin{itemize}
 \item \textbf{The two base models}, Kabbalah and Tibetan Buddhism, treated in the most detail, because each
 built a full and explicit system that practitioners have refined for many centuries.
 \item \textbf{The six living esoteric traditions}, Hinduism, Christianity, Sufism, Sikhism,
 Taoism, and Hermeticism, each rich and living, treated in solid but lighter depth.
 \item \textbf{The six ancient cosmologies}, Greek, Roman, Norse, Egyptian, Mesopotamian, and Chinese, more
 descriptive and anthropomorphic, told with their gods and their imagery.
 \item \textbf{Eighteen indigenous and less-documented cosmologies}, each kept brief, a creation story and a
 few key motifs, to show the global reach of the same patterns.
 \item \textbf{A synthesis}, drawing the common threads together and proposing the one model behind the
 many.
 \item \textbf{An appendix}, giving for each tradition a mapping from plain modern terms to its native
 vocabulary, and its creation story retold once more in its own language.
\end{itemize}

Every chapter opens the same way, with the creation story in plain modern English and no special terms, so the
shape comes first and the vocabulary second. This is deliberate. The native words are powerful and precise,
but they can also wall a tradition off and hide how much it shares with its neighbors. By telling the story
twice, once bare and once clothed, we can see both the shared structure and the unique detail.
